\section{Introduction}

Differentiable Logic Gate Networks (DLGNs) provide an alternative path toward
efficient machine learning at the edge by learning directly with the primitive operations
used in digital hardware. Rather than relying on multiply-accumulate
operations, a DLGN is composed of Boolean gates. During training, these discrete operations are
represented through differentiable relaxations, allowing the network to be
optimized with gradient descent. After training, each relaxed gate is hardened
to a Boolean function for efficient inference \cite{petersen2022}. This direct
correspondence between the learned model and Boolean computation makes DLGNs
particularly attractive for resource-constrained embedded and edge and systems.

However, a fundamental characteristic of DLGNs is their sparse
connectivity. Each conventional DLGN gate has only two inputs. These two predecessors are selected before training and remain
fixed, while optimization determines which Boolean function the gate should
implement \cite{petersen2022}. Therefore, training can change \emph{how} two
signals are combined, but not \emph{which} signals arrive at a gate. To that end,  if two
informative signals are never connected through the available graph, learning
the gate functions alone cannot create that missing edge. At a fixed depth and
width, connectivity consequently determines the combinations of information
that can be formed throughout the network and becomes part of the effective
architectural capacity of the DLGN.

A way to address this issue is to consider connectivity as an independent
design dimension~\cite{buehrer2026} and optimize the predecessors themselves during
training \cite{mommen2025,fojcik2026}. This type of learned connectivity can recover
connections that were absent from the original graph and can therefore improve
accuracy. However, such methods require significantly increased training time as each gate must consider multiple candidate signals,
requiring additional routing state and candidate activations even though the
final hardened network may retain only a small number of predecessors. This
creates a practical design choice between inexpensive but arbitrary
\emph{fixed random connectivity} and more flexible but substantially heavier
\emph{learned connectivity}.


% The connectivity question becomes even more important for vision workloads.
% Early DLGNs operate on flattened image representations, where random
% connections do not explicitly preserve the strong spatial and channel
% relationships present in images \cite{petersen2022}. Convolutional DLGNs
% address this limitation by introducing shared logic-gate trees that operate
% over local receptive fields, enabling logic networks to capture spatial
% patterns and scale to substantially more challenging vision tasks
% \cite{petersen2024}. Yet convolution does not remove the connectivity design
% problem; instead, it changes its form. A convolutional logic tree must select
% both \emph{where} to sample within its receptive field and \emph{which input
% channels} provide those samples. The spatial pattern is reused across image
% locations through weight sharing, making arbitrary changes to this structure
% undesirable. Thus, a useful connectivity method should respect the
% convolutional spatial sampler while improving how channels are paired. A
% shared connectivity principle must therefore accommodate two distinct
% representations: individual input coordinates in dense DLGNs and shared
% channel groups in convolutional DLGNs.

% To demonstrate the need for connectivity design, we consider a
% motivational example on a dense CIFAR-10 DLGN. We keep the deeper
% connectivity unchanged and randomize only the first-layer connections
% under a degree-preserving constraint. Hence, every predecessor is used
% exactly the same number of times before and after randomization; only
% the identities of the signals paired at each gate are changed. Despite
% preserving the same fan-out, hardened validation accuracy decreases by
% $4.49$ percentage points. This result reveals an important limitation of
% degree-balanced connectivity: \emph{how often a signal is used does not
% determine which information is brought together}. However, explicitly
% selecting favorable pairs is itself challenging to scale. Given $N$
% candidate predecessors, a two-input gate can already choose from
% $\mathcal{O}(N^2)$ possible pairs, and this choice is repeated across
% thousands or millions of gates and multiple network layers. Moreover,
% while input signals have explicit image semantics, intermediate signals
% progressively lose this direct interpretation, making a manually defined
% pairing rule increasingly difficult to justify at deeper layers. Fixed
% random connectivity avoids this search but ignores signal relationships,
% whereas learned connectivity explores the enlarged space by maintaining
% additional routing state and candidate activations during training.
% Therefore, a key challenge is to construct informative connections
% \emph{before training} using principles that remain scalable as the
% network grows and do not introduce the training overhead of a separate
% routing problem.


To isolate the effect of \emph{pair identity}, we consider a four-layer
dense CIFAR-10 DLGN with 512K gates. We compare its structured first-layer
graph against a randomized control derived from the same graph. The control
preserves exactly how often every encoded input appears in each of the two
gate-input positions, but permutes these assignments across gates. Thus, the
two networks have identical dimensions and predecessor-usage statistics and
differ only in which inputs are paired. All deeper connections remain
unchanged, and both networks use the same data split, paired initialization,
and training configuration.
\begin{table}[t]
\centering
\caption{Effect of first-layer pair identity on a four-layer, 512K-gate
dense CIFAR-10 DLGN. The randomized control preserves, for every encoded
input, its number of occurrences in each gate-input slot; only the
pairings across gates are changed. All deeper connections and training
settings are identical. Values report best hardened validation accuracy
(\%).}
\label{tab:motivation}
\small
\setlength{\tabcolsep}{4pt}
\begin{tabular}{c|cc|c}
\hline
\textbf{Seed} &
\textbf{Structured Pairs} &
\textbf{Randomized Pairs} &
$\mathbf{\Delta}$ \textbf{(pp)} \\
\hline
0 & 58.86 & 54.26 & +4.60 \\
1 & 59.48 & 55.26 & +4.22 \\
2 & 59.06 & 54.40 & +4.66 \\
\hline
\textbf{Mean $\pm$ SD}
& $\mathbf{59.13\pm0.32}$
& $54.64\pm0.54$
& $\mathbf{+4.49\pm0.24}$ \\
\hline
\end{tabular}
\end{table}
As shown in Table~\ref{tab:motivation}, randomizing only the pair identities
reduces hardened validation accuracy in all three runs, with a mean drop of
$4.59\%$. Hence, balanced predecessor usage alone does not
determine connectivity quality: the specific signals combined by each gate
also matter. Determining favorable pairs, however, quickly becomes
intractable. With $N$ available predecessors, each two-input gate admits
$N(N-1)/2$ possible unordered pairs, and this decision is repeated across
many gates and layers. Furthermore, image inputs provide explicit spatial,
channel, and encoding relationships that can guide connectivity, whereas
intermediate signals generally lack the same direct interpretation. This
creates the need for a scalable construction that exploits known signal
relationships where available and organizes subsequent interactions without
introducing a separate trainable routing problem.

In this paper, we propose \textbf{LogicConnect}, an offline connectivity
construction that introduces structure into the DLGN graph before training.
LogicConnect first exploits signal identity where it is most explicit by
forming semantic first-layer pairs according to relationships among the
encoded inputs rather than arbitrary indices. In subsequent layers, it
organizes predecessor interactions through regular multiscale matchings,
selecting complete matching stages primarily to balance predecessor usage
and using ancestry overlap to distinguish alternatives with equivalent
degree balance. 

% This leads to the central question investigated in this work:
% \emph{how much accuracy can a DLGN recover through structured offline
% connectivity, under the same declared architecture, before the additional
% training cost of learned routing becomes necessary?}
% Our main contributions are:

% \begin{itemize}

%     \item We propose \textbf{LogicConnect}, a shared fixed-connectivity
%     construction for dense and convolutional DLGNs. It combines semantic
%     input pairing with regular multiscale predecessor matchings, while
%     preserving the conventional fixed-index representation and introducing no
%     trainable routing parameters.

%     \item We quantify the resulting \textbf{accuracy--training-resource
%     trade-off}. Across three matched dense CIFAR-10 network sizes,
%     LogicConnect improves hardened test accuracy over fixed random
%     connectivity by $3.18$--$4.59$ percentage points across three paired
%     seeds. Compared with the evaluated Top-32 learned-routing configuration,
%     it uses $93.7$--$94.2\%$ less peak GPU memory while remaining within
%     $1.84$--$5.30$ accuracy points. We further evaluate the same connectivity
%     principle in convolutional DLGNs; the held-out single-seed results are
%     positive, while replicated validation results remain inconclusive.

%     \item We perform \textbf{structural attribution} to identify where the
%     gains originate. Degree-preserving experiments separate predecessor
%     identity from predecessor frequency and show the strongest measured dense
%     effect at the first layer. Additional experiments examine the incremental
%     effect of deeper connectivity, simpler construction rules, alternative
%     gate training, and distribution shift, characterizing both the benefits
%     and the limits of structured fixed connectivity.

%     \item LogicConnect neither introduces a new inference operator nor
% increases the declared depth, width, or gate budget; instead, it moves
% connectivity design entirely offline, providing a more informative fixed
% graph without creating an additional trainable routing problem.
% \item differences between desne and convolutional

% \end{itemize}   



\begin{comment}

\section{Introduction}

A differentiable logic gate network (DLGN) is constrained not only by the Boolean functions learned at its gates, but also by the signals that those gates are allowed to access. Each gate has only two predecessors; once these predecessors are fixed, gradient-based training can optimize how the two signals are combined, but it cannot introduce a connection that is absent from the graph. Consequently, the connectivity selected before training determines which signal combinations can be formed at a given depth and width, and therefore constitutes an architectural design choice rather than merely an implementation detail. This distinction is particularly important for DLGNs, where the original formulation learns a distribution over Boolean functions while keeping randomly initialized connections fixed throughout training \cite{petersen2022}.

This constraint creates a practical design choice. Fixed random connectivity is inexpensive: each gate stores only its two predecessor indices and training is restricted to the gate-function parameters \cite{petersen2022,petersen2024}. More recent approaches instead make connectivity trainable, allowing a gate to select among alternative predecessors during optimization \cite{mommen2025,fojcik2026}. Such learned routing can substantially improve the usefulness of a given gate budget, but it also requires representing and processing candidate connections during training. The final hardened circuit may still contain only two inputs per gate, yet obtaining those connections can require considerably more training state and intermediate activations. This motivates the central question of this work: \emph{how much accuracy can be recovered by designing the fixed graph before training, without increasing the declared gate budget or introducing a second trainable routing problem?}

This question is complementary to several recent directions in DLGN research. Convolutional DLGNs extend the underlying architecture through shared logic-gate trees and spatial processing \cite{petersen2024}; other work improves the continuous-to-discrete gate parameterization and reduces the discretization gap \cite{yousefi2025}; learned-connectivity methods optimize the network connectome itself \cite{mommen2025,fojcik2026}; and post-training techniques simplify an already trained circuit to improve its hardware footprint \cite{lee2026}. In contrast, we focus specifically on the \emph{fixed connection graph}: the graph is constructed once, before gate-function training, and subsequently uses the same fixed-index representation as a conventional DLGN. Connectivity design is therefore orthogonal to improvements in gate relaxation and post-training circuit simplification, and offers a different accuracy--training-resource operating point from learned routing.

Importantly, useful connectivity cannot be reduced to balancing predecessor degree. The native dense random construction already distributes predecessor usage broadly, yet it does not control \emph{which} signals are paired. Two graphs can therefore exhibit the same predecessor degrees while exposing their gates to very different combinations of information. We observe this directly in a degree-preserving ablation on the dense CIFAR-10 medium configuration: randomizing only the first-layer predecessor identities while preserving the degree of every predecessor in each input slot reduces hardened validation accuracy by $4.49$ percentage points when the deeper graph remains structured. Across three paired seeds, the corresponding $95\%$ confidence interval is $[3.90,5.09]$. Thus, predecessor identity has a measurable effect beyond fan-out balance. This observation motivates treating connectivity as a structured architectural variable and, in particular, exploiting information that is discarded by an index-only random assignment.

In this work, we introduce \textbf{LogicConnect}, an offline fixed-connectivity construction for both dense and convolutional DLGNs. LogicConnect first forms \emph{semantic input pairs} using the source structure of image features, so that the first logic layer combines inputs according to their channel, spatial, and binary-encoding relationships rather than arbitrary indices. For subsequent layers, it constructs regular \emph{multiscale predecessor matchings} that expose gates to different interaction scales while maintaining balanced predecessor usage. Whole matching stages are selected using degree balance as the primary criterion and ancestry overlap to distinguish equally balanced alternatives. Once constructed, the graph remains unchanged and ordinary DLGN training learns only the gate functions.

The same construction is adapted to convolutional DLGNs without replacing their convolutional execution model. LogicConnect changes the \emph{channel pairs} presented to the shared logic-gate trees while retaining the existing spatial sampler and weight-sharing mechanism of the convolutional architecture \cite{petersen2024}. The subsequent classifier uses the same fixed structured-connectivity principle as the dense network. Hence, LogicConnect does not introduce a new routing operator at inference and does not enlarge the declared network width, depth, or gate count; it exchanges a one-time offline graph-construction step for a more deliberate set of signal interactions during training.

Overall, our main contributions are:

\begin{itemize}

    \item We introduce \textbf{LogicConnect}, a shared fixed-connectivity construction for dense and convolutional DLGNs. It combines semantic input pairing with balanced multiscale predecessor matchings while retaining the conventional two-index-per-gate representation and requiring no trainable routing parameters.

    \item We quantify the resulting \textbf{accuracy--training-resource trade-off}. Under matched dense CIFAR-10 architectures, LogicConnect improves hardened test accuracy over fixed random connectivity by $3.18$--$4.59$ percentage points across three evaluated network sizes and three paired seeds. Relative to the evaluated learned Top-32 routing comparator, LogicConnect requires $93.7$--$94.2\%$ less peak GPU memory while remaining within $1.84$--$5.30$ accuracy points. The convolutional implementation further demonstrates applicability of the shared construction: held-out single-seed results are positive, while replicated validation evidence remains statistically inconclusive.

    \item We perform \textbf{structural attribution} of the measured gains. Degree-preserving ablations separate predecessor identity from predecessor frequency and identify the first layer as the strongest measured source of the dense improvement. Additional experiments examine the incremental effect of deeper connectivity, simpler construction rules, alternative gate training, and distribution shift, exposing both the benefits and the limits of structured fixed connectivity.

\end{itemize}
\end{comment}